\Article{14}
Einen Umstand unserer allgemeinen Lösung müssen wir noch etwas ausführlicher betrachten. In Artikel 5 haben wir gezeigt, dass es immer zwei Lösungen gibt. Bei der einen muss \(P+iQ\) eine Funktion von \(p+iq\) und \(P-iQ\) eine Funktion von \(p-iq\) sein. Bei der anderen ist \(P+iQ\) eine Funktion von \(p-iq\) und \(P-iQ\) eine Funktion von \(p+iq\). Wir wollen nun zeigen, dass bei der einen Lösung die Teile der Abbildung immer ähnlich angeordnet sind wie im Urbild, während sie bei der anderen seitenverkehrt liegen. Zugleich wollen wir das Kriterium angeben, nach dem sich das im Voraus unterscheiden lässt.

Zunächst bemerken wir, dass von vollkommener oder seitenverkehrter Ähnlichkeit nur die Rede sein kann, wenn man an jeder der beiden Flächen zwei Seiten unterscheidet und eine als obere, die andere als untere betrachtet. Da diese Wahl an sich etwas willkürlich ist, unterscheiden sich die beiden Lösungen im Wesentlichen nicht. Eine seitenverkehrte Ähnlichkeit wird zur vollkommenen, sobald man bei einer der Flächen die zuvor als obere betrachtete Seite zur unteren macht. In unserer Lösung konnte diese Unterscheidung daher keine Rolle spielen, denn die Flächen waren allein durch die Koordinaten ihrer Punkte bestimmt. Will man die Unterscheidung berücksichtigen, muss man die Flächen zunächst auf eine andere Weise festlegen, die sie einschließt. Zu diesem Zweck nehmen wir an, die erste Fläche sei durch die Gleichung \(\psi=0\) bestimmt, wobei \(\psi\) eine gegebene eindeutige Funktion von \(x,y,z\) ist. In allen Punkten der Fläche verschwindet also der
\SourceLeaf{217}
Wert von \(\psi\), und in allen Punkten des Raumes außerhalb der Fläche ist er von null verschieden. Beim Durchgang durch die Fläche geht der Wert von \(\psi\) also, wenigstens im Allgemeinen, vom Positiven ins Negative über und beim Durchgang in der Gegenrichtung vom Negativen ins Positive. Mit anderen Worten, der Wert von \(\psi\) ist auf der einen Seite positiv und auf der anderen negativ. Die erste betrachten wir als obere, die andere als untere Seite. Bei der zweiten Fläche verfahren wir genauso. Sie ist durch die Gleichung \(\Psi=0\) bestimmt, wobei \(\Psi\) eine gegebene eindeutige Funktion der Koordinaten \(X,Y,Z\) ist. Die Differentiation ergebe ferner
\[
\begin{aligned}
\mathrm{d}\psi&=e\,\mathrm{d}x+g\,\mathrm{d}y+h\,\mathrm{d}z\\
\mathrm{d}\Psi&=E\,\mathrm{d}X+G\,\mathrm{d}Y+H\,\mathrm{d}Z
\end{aligned}
\]
\noindent wobei \(e,g,h\) Funktionen von \(x,y,z\) und \(E,G,H\) Funktionen von \(X,Y,Z\) sind.

Die Überlegungen, mit denen wir zum angestrebten Ziel gelangen, sind an sich nicht schwierig, aber etwas ungewöhnlich. Wir wollen sie daher so klar wie möglich darstellen. Zwischen den beiden einander entsprechenden Darstellungen auf den Flächen mit den Gleichungen \(\psi=0\) und \(\Psi=0\) nehmen wir sechs Zwischendarstellungen in der Ebene an. Damit betrachten wir acht verschiedene Darstellungen.
\begin{center}
\begingroup\fontsize{8}{10}\selectfont\setlength{\tabcolsep}{2pt}
\begin{tabular}{@{}lll@{}}
\multicolumn{3}{r}{\parbox{1.55in}{\raggedright als entsprechend gelten die Punkte mit den jeweiligen Koordinaten \(=\)}}\\[2pt]
1\textsuperscript{s}&das Urbild auf der Fläche mit der Gleichung \(\psi=0\) . .&\(x,\ y,\ z\)\\
2\textsuperscript{s}&Darstellung in der Ebene . . . . . . . . . . . . . . . . . . . . .&\(x,\ y,\ 0\)\\
3\textsuperscript{s}&\quad „\quad „\quad „\quad „\quad . . . . . . . . . . . . . . . . . . . . .&\(t,\ u,\ 0\)\\
4\textsuperscript{s}&\quad „\quad „\quad „\quad „\quad . . . . . . . . . . . . . . . . . . . . .&\(p,\ q,\ 0\)\\
5\textsuperscript{s}&\quad „\quad „\quad „\quad „\quad . . . . . . . . . . . . . . . . . . . . .&\(P,\ Q,\ 0\)\\
6\textsuperscript{s}&\quad „\quad „\quad „\quad „\quad . . . . . . . . . . . . . . . . . . . . .&\(T,\ U,\ 0\)\\
7\textsuperscript{s}&\quad „\quad „\quad „\quad „\quad . . . . . . . . . . . . . . . . . . . . .&\(X,\ Y,\ 0\)\\
8\textsuperscript{s}&Abbildung auf der Fläche mit der Gleichung \(\Psi=0\) . .&\(X,\ Y,\ Z\)
\end{tabular}\endgroup
\end{center}
\par\indent Wir wollen nun diese Darstellungen allein hinsichtlich der gegenseitigen \textit{Lage} der unendlich kleinen Linienelemente vergleichen %
\SourceLeaf{218}%
und das Größenverhältnis ganz beiseitelassen. Zwei Darstellungen gelten als ähnlich angeordnet, wenn von zwei Linienelementen, die von einem Punkt ausgehen, dem in der einen Darstellung rechts liegenden Element auch in der anderen das rechts liegende entspricht. Andernfalls nennen wir ihre Lage seitenverkehrt. In der Ebene bei Nr. 2---7 gilt stets die Seite mit den positiven Werten der dritten Koordinate als obere. Bei der ersten und der letzten Fläche hängt die Unterscheidung zwischen oberer und unterer Seite dagegen allein davon ab, ob \(\psi\) und \(\Psi\) positiv oder negativ sind, wie oben festgelegt.

Zunächst ist klar, dass die Darstellung 2 an jeder Stelle der ersten Fläche ähnlich angeordnet ist wie Darstellung 1, an der man bei unverändertem \(x\) und \(y\) durch einen positiven Zuwachs von \(z\) auf die obere Seite gelangt. Das ist offenbar überall dort der Fall, wo \(h\) positiv ist. Bei negativem \(h\) tritt das Gegenteil ein, und die Darstellungen sind seitenverkehrt angeordnet.

Entsprechend sind die Darstellungen 7 und 8 ähnlich oder seitenverkehrt angeordnet, je nachdem, ob \(H\) positiv oder negativ ist.

Zum Vergleich der Darstellungen 2 und 3 sei in der ersten \(\mathrm{d}s\) die Länge einer unendlich kleinen Linie vom Punkt mit den Koordinaten \(x,y\) zu einem anderen Punkt mit den Koordinaten \(x+\mathrm{d}x,y+\mathrm{d}y\). Ihre Neigung gegen die Abszissenachse sei \(l\) und wachse in dem Sinn, in dem man von der \(x\)-Achse zur \(y\)-Achse übergeht. Dann gilt
\[
\mathrm{d}x=\mathrm{d}s\mathord{.}\cos l,\quad \mathrm{d}y=\mathrm{d}s\mathord{.}\sin l
\]
\noindent In Darstellung 3 sei \(\mathrm{d}\sigma\) die Länge der Linie, die \(\mathrm{d}s\) entspricht. Ihre Neigung gegen die Abszissenachse sei im selben Sinn wie zuvor \(\lambda\), sodass
\[
\mathrm{d}t=\mathrm{d}\sigma\mathord{.}\cos\lambda,\quad \mathrm{d}u=\mathrm{d}\sigma\mathord{.}\sin\lambda
\]
\noindent In der Schreibweise aus Artikel 4 gilt also
\[
\begin{aligned}
\mathrm{d}s\mathord{.}\cos l&=\mathrm{d}\sigma(a\cos\lambda+a'\sin\lambda)\\
\mathrm{d}s\mathord{.}\sin l&=\mathrm{d}\sigma(b\cos\lambda+b'\sin\lambda)
\end{aligned}
\]
\noindent und folglich
\[
\mathrm{tang}\,l=\frac{b\cos\lambda+b'\sin\lambda}{a\cos\lambda+a'\sin\lambda}
\]
